\section{Background and Related Work}\label{sec:related}

\subsection{Decentralized Networks}
Since Chaum's seminal work on untraceable email in 1981~\cite{chaum-mix}, there has been a great amount of research related to mixnet design~\cite{piotrowska2017loopix, van2015vuvuzela, kwon2020xrd, lazar2018karaoke, cottrell1995mixmaster, alexopoulos2017MCMIX, chaum2016cmix, chaum-mix, danezis2003mixminion}. 
A mixnet is a network of nodes that delay and shuffle packets as they traverse a sequence of hops before reaching their final destination, preventing adversaries from correlating input and output traffic.

Mixnets have been known to resist powerful adversaries typically called global passive adversaries whose goal is to deanonymize clients and trace message through traffic analysis techniques. 
Mixnets achieve this by (i) mixing messages inside the mixnodes, thereby hiding the connection between incoming and outgoing packets and (ii) the use of a cryptographic packet format that prevents adversaries from tracing messages based on packet size or format. 
Messages are first divided into fixed-size packets to prevent size-based correlation attacks. 
Each packet is then encrypted in multiple layers, together with the necessary routing information. 
This guarantees that every mixnode learns only its immediate predecessor and successor along the path.

The Lightning Network~\cite{poon2016bitcoin} is a payment-channel that allow two users to send off-chain transactions without continuously interacting with the Bitcoin blockchain. 
In this design, even users who do not have a direct channel may still transact by routing payments through intermediate peers. 
When a node wants to pay another node without having a direct channel, the transaction originator first computes the shortest path to the recipient and then encrypts the routing information multiple times, where each cryptographic layer corresponds to a hop along the path.
Although the Lightning Network was initially introduced to address Bitcoin's scalability, it also addresses the privacy issues introduced by peer-to-peer networks~\cite{sharma2022anonymity}. 
In particular adversaries should not be able to identify the sender or receiver of a transaction. 
Additionally an intermediate node should not know whether their predecessor is the transaction's originator or whether their successor is the transaction's recipient. 

These privacy guarantees, both in mixnets and in the Lightning Network, are achieved using the Sphinx packet format~\cite{sphinx}.
\begin{comment}
\begin{figure}
	\hspace{-9mm}
	\resizebox{1.085\linewidth}{!}{
\newcommand{\figonetextcol}[0]{\color{black}}

\tikzset{  
    innerLink/.style={-, densely dotted, shorten <= 1pt, shorten >= 2pt},  
    outerLink/.style={->, shorten <= 2pt, shorten >= 4pt, >=Latex},
    packet/.style={circle, inner sep=0.2pt, fill=#1},
    A/.style={packet=\colorA!40, font={#1}},
    B/.style={packet=\colorB!50, font={#1}},
    A1/.style={A=1}, A2/.style={A=2}, A3/.style={A=3}, A4/.style={A=4}, A5/.style={A=5}, A6/.style={A=6}, A7/.style={A=7}, A8/.style={A=8}, A9/.style={A=9},
    B1/.style={B=1}, B2/.style={B=2}, B3/.style={B=3}, B4/.style={B=4}, B5/.style={B=5}, B6/.style={B=6}, B7/.style={B=7}, B8/.style={B=8}, B9/.style={B=9},
}

\newcommand{\countitems}[1]{%
  \def\size{0}%
  \foreach \x in {#1} {%
    \pgfmathparse{\size+1}%
    \global\let\size=\pgfmathresult%
  }%
}

\newcommand{\outerLink}[3]{
  % size (nbr of packets)
  \countitems{#3} 
  % If odd nbr of packet, slight shift to not have a packet on an intersection
  \pgfmathparse{mod(\size,2)==1 ? \size-0.5 : \size}%
  \xdef\size{\pgfmathresult}%

  % define line
  \coordinate (start) at (#1.east);
  \coordinate (stop) at (#2.west);
  \draw[outerLink] (start) -- (stop);
  % put packet on the line
  \foreach \packet [count=\i from 1] in {#3} {
    \pgfmathsetmacro{\t}{1 - \i/(\size+1)}%
    \path let
      \p1 = (start),
      \p2 = (stop),
      \p3 = ($(\p1)!\t!(\p2)$)
    in
      node[\packet] at (\p3) {};
  }
}


\newcommand{\mixnode}[2]{%
  \countitems{#1}%

  \begin{tikzpicture}
    \def\d{0.8}
    \draw[thick] (0,0) circle (\d);

    % Connect according to mapping
    \pgfmathsetmacro{\xoffset}{(1.14*\d - \size/10)*\d}
    \foreach \r [count=\l from 1] in {#1} {
      \pgfmathsetmacro{\yposL}{\d - \l/(\size+1)*2*\d}
      \coordinate (left\l) at (-\xoffset, \yposL);
      \pgfmathsetmacro{\yposR}{\d - \r/(\size+1)*2*\d}
      \coordinate (right\l) at (\xoffset, \yposR);
    }
    % Draw dots
    \foreach \packet [count=\i from 1] in {#2} {
      \draw[innerLink] (left\i) -- (right\i);
      \node[\packet] at (left\i) {};
      \node[\packet] at (right\i) {};
    }
  \end{tikzpicture}
}

\newcommand{\client}[1]{
  \begin{tikzpicture}[scale=0.4]
    % Monitor screen
    \draw[fill=#1!20, draw=black, rounded corners=4pt] (0,0) rectangle (4,2.5);
    % Inner bezel
    \draw[fill=white, draw=none, rounded corners=3pt] (0.2,0.2) rectangle (3.8,2.3);
    % Stand
    \draw[fill=#1!40, draw=black] (1.75,-0.7) rectangle (2.25,0);
    % Base
    \draw[fill=#1!40, draw=black, rounded corners=2pt] (0.8,-1) rectangle (3.2,-0.7);
  \end{tikzpicture}
}

\newcommand{\server}[1]{
  \begin{tikzpicture}[scale=0.7]
    \foreach \y in {0, -0.6, -1.2} {
        \draw[fill=#1!20] (0,\y) rectangle (2,\y-0.5); % rectangle
        \draw (0.2,\y-0.25) circle (0.05); % power/status LED
    }
  \end{tikzpicture}
}

\newcommand{\colorA}{blue}
\newcommand{\colorB}{orange}

\newcommand{\clientA}{\client{\colorA}}
\newcommand{\clientB}{\client{\colorB}}
\newcommand{\serverA}{\server{\colorA}}
\newcommand{\serverB}{\server{\colorB}}


\begin{tikzpicture}
  \def\dx{4}
  \def\dy{-1}

  % GENERAL ARCHITECTURE (define nodes)
  \node (clientA) at (0,-\dy) {\clientA};
  \node (clientB) at (0,\dy) {\clientB};

  \node (serverA) at (4*\dx,-\dy) {\serverA};
  \node (serverB) at (4*\dx,\dy) {\serverB};

  % PACKETS - links and shuffles (python generated)
\node (mix11) at (1*\dx,-\dy) {\mixnode{3,1,2}{A2,A1,B3}};
\node (mix12) at (1*\dx, \dy) {\mixnode{3,2,1}{B2,B1,A3}};
\node (mix21) at (2*\dx,-\dy) {\mixnode{3,1,2}{A3,A2,B3}};
\node (mix22) at (2*\dx, \dy) {\mixnode{2,3,1}{B2,B1,A1}};
\node (mix31) at (3*\dx,-\dy) {\mixnode{3,1,4,2}{B2,A2,A1,B3}};
\node (mix32) at (3*\dx, \dy) {\mixnode{2,1}{A3,B1}};

\outerLink{clientA}{mix11}{A1,A2};
\outerLink{clientA}{mix12}{A3};
\outerLink{clientB}{mix11}{B3};
\outerLink{clientB}{mix12}{B1,B2};

\outerLink{mix11}{mix21}{A2,B3};
\outerLink{mix11}{mix22}{A1};
\outerLink{mix12}{mix21}{A3};
\outerLink{mix12}{mix22}{B2,B1};

\outerLink{mix21}{mix31}{A2,B3};
\outerLink{mix21}{mix32}{A3};
\outerLink{mix22}{mix31}{B2,A1};
\outerLink{mix22}{mix32}{B1};

\outerLink{mix31}{serverA}{A2,A1};
\outerLink{mix31}{serverB}{B2,B3};
\outerLink{mix32}{serverA}{A3};
\outerLink{mix32}{serverB}{B1};

  % ANNOTATIONS
  \node at (0*\dx, -2.2*\dy) {\figonetextcol Clients};
  \draw[loosely dotted, fill=gray!30] (0.23*\dx, -2.2*\dy) -- (0.23*\dx, +2*\dy) {};
  \node at (0.5*\dx, -2.2*\dy) {\figonetextcol $t_0$};
  \draw[loosely dotted, fill=gray!30] (0.77*\dx, -2.2*\dy) -- (0.77*\dx, +2*\dy) {};
  
  \node at (1*\dx, -2.2*\dy) {\figonetextcol $t_1$};
  \draw[loosely dotted, fill=gray!30] (1.23*\dx, -2.2*\dy) -- (1.23*\dx, +2*\dy) {};
  \node at (1.5*\dx, -2.2*\dy) {\figonetextcol $t_2$};
  \draw[loosely dotted, fill=gray!30] (1.77*\dx, -2.2*\dy) -- (1.77*\dx, +2*\dy) {};
  
  \node at (2*\dx, -2.2*\dy) {\figonetextcol $t_3$};
  \draw[loosely dotted, fill=gray!30] (2.23*\dx, -2.2*\dy) -- (2.23*\dx, +2*\dy) {};
  \node at (2.5*\dx, -2.2*\dy) {\figonetextcol $t_4$};
  \draw[loosely dotted, fill=gray!30] (2.77*\dx, -2.2*\dy) -- (2.77*\dx, +2*\dy) {};
  
  \node at (3*\dx, -2.2*\dy) {\figonetextcol $t_5$};
  \draw[loosely dotted, fill=gray!30] (3.23*\dx, -2.2*\dy) -- (3.23*\dx, +2*\dy) {};
  \node at (3.5*\dx, -2.2*\dy) {\figonetextcol $t_6$};
  \draw[loosely dotted, fill=gray!30] (3.77*\dx, -2.2*\dy) -- (3.77*\dx, +2*\dy) {};
  \node at (4*\dx, -2.2*\dy) {\figonetextcol Services};

\end{tikzpicture}
}
	\caption{
		Step-by-step tracing of two clients communicating with services through a 3-layer mixnet with two mixnodes (large circles) at each layer.
	}
	\label{fig:mixnet}
\end{figure}\todo{we can potentially delete this picture}

\end{comment}


\subsection{Sphinx Packet Format}
% Sphinx
The \emph{Sphinx} packet format provides three essential security properties: per-hop confidentiality, per-hop integrity protection, and strong unlinkability.

First, Sphinx ensures \emph{per-hop confidentiality} by encrypting routing information in layers, similar to classical onion routing.  
Each mixnode learns only its immediate predecessor and successor, but gains no knowledge about the overall path, the sender, or the final destination. Second, Sphinx provides \emph{per-hop integrity protection} through a cryptographic tag embedded in each layer of the header. These tags prevent adversaries from tampering with routing information.  Any modification results in an integrity-tag verification failure at the routing node, causing the packet to be discarded.
This ensures that adversaries cannot tamper with routes or perform active attacks such as header tagging in order to target specific messages or clients. Third, Sphinx offers \emph{strong unlinkability} by ensuring that packets entering and leaving a mixnode are computationally indistinguishable.  
This is achieved by using fixed-size headers and payloads and by re-randomizing cryptographic elements at every hop.
Each hop transforms the packet into a fresh, uniformly distributed header, preventing an adversary from linking an incoming packet to its outgoing counterpart. These guarantees have made Sphinx the de facto standard for privacy-preserving routing in modern decentralized systems such as the Lightning Network~\cite{poon2016bitcoin} and mixnet-based architectures like Nym~\cite{nym} and Katzenpost~\cite{infeld2025echomix}.

However, in all these systems, the Sphinx packet format must be fully constructed by the client, which introduces vulnerabilities when malicious clients deviate from the prescribed routing algorithm. 
Routing algorithms in decentralized networks are chosen carefully balancing latency with strong privacy guarantees. Preventing malicious clients from deviating from the routing algorithm, without relying on a central authority and while preserving unlinkability between packets and clients has been a challenging open problem. A malicious client can intentionally deviate from the routing protocol to launch attacks, either against the functioning of the system, against specific clients, or against the overall privacy provided by the system. For example, in the Lightning Network, a malicious client can route payments through long paths specifically to launch jamming attacks, freeze funds, or infer activity patterns of nodes through response times~\cite{Rohrer2019paymentchannels,kappos2021empirical}. In Nym, nodes should be chosen uniformly across layers to maximize anonymity~\cite{benguirat2022mixnet}; therefore, an attacker can bias route selection either to reduce the overall anonymity of the system or to target specific clients and messages.

\begin{comment}
In addition, enforcing policies such as accountability and legitimacy verification (exit policy) remains inherently challenging in decentralized networks\cite{oram2001peer,diaz2007accountable, naylor2014balancing,mirkovic2008building, backes2014backref, xu2012accountable}.
Since Sphinx headers are constructed to reveal minimal information to mixnodes, networks cannot reliably verify client credentials, enforce usage policies, or apply access control before packets reach the exit node.
This trade-off between strong privacy and limited accountability, coupled with the risks posed by untrusted or misbehaving clients, presents practical obstacles for the widespread deployment of mixnet-based systems.
\end{comment}
